\documentclass[a4paper, 12pt]{article}
\usepackage[utf8]{inputenc}
\usepackage[spanish]{babel} % Cambia a 'english' si prefieres inglés
\usepackage{graphicx}
\usepackage{float}
\usepackage{amsmath}
\usepackage{booktabs}
\usepackage{array}
\usepackage{caption}
\usepackage{subcaption}
\usepackage{circuitikz} % Para diagramas de circuitos eléctricos
\usepackage{hyperref}
\usepackage{geometry}
\usepackage{listings}
\usepackage{xcolor}
\usepackage{siunitx} % Paquete para unidades del SI

% Configuración de márgenes
\geometry{
    a4paper,
    left=2.5cm,
    right=2.5cm,
    top=2.5cm,
    bottom=2.5cm
}

% Configuración de hyperref
\hypersetup{
    colorlinks=true,
    linkcolor=blue,
    filecolor=magenta,      
    urlcolor=cyan,
    pdftitle={Documento Técnico de Circuitos},
    pdfpagemode=FullScreen,
}

% Configuración para listados de código (opcional)
\lstset{
    language=Python,
    basicstyle=\ttfamily\small,
    keywordstyle=\color{blue},
    commentstyle=\color{gray},
    numbers=left,
    numberstyle=\tiny\color{gray},
    frame=single,
    breaklines=true,                % Permite que las líneas largas se corten
    captionpos=b,                   % Posición del caption (abajo)
    literate=
      {á}{{\'a}}1 {é}{{\'e}}1 {í}{{\'i}}1 {ó}{{\'o}}1 {ú}{{\'u}}1
      {Á}{{\'A}}1 {É}{{\'E}}1 {Í}{{\'I}}1 {Ó}{{\'O}}1 {Ú}{{\'U}}1
      {ñ}{{\~n}}1 {Ñ}{{\~N}}1
      {¿}{{?`}}1 {¡}{{!`}}1
}

% Título del documento
\title{\textbf{Plantilla para Crear Documentos con Circuitos Eléctricos}}
\author{Prof. César Augusto Rodríguez}
\date{\today}

\begin{document}

\maketitle
\thispagestyle{empty}

\begin{abstract}
Este documento presenta una plantilla para crear informes técnicos que incluyen esquemas de circuitos eléctricos, imágenes en formato PNG y tablas de datos. La plantilla utiliza los paquetes \texttt{circuitikz} para diagramas de circuitos y \texttt{graphicx} para imágenes.
\end{abstract}

\tableofcontents
\newpage

\section{Introducción}
Esta plantilla está diseñada para facilitar la creación de documentos técnicos relacionados con circuitos eléctricos. Incluye ejemplos de cómo insertar:
\begin{itemize}
    \item Diagramas de circuitos usando \texttt{circuitikz}
    \item Imágenes en formato PNG
    \item Tablas de datos organizadas
    \item Fórmulas matemáticas
\end{itemize}

\section{Circuitos Eléctricos}

\subsection{Circuito Básico RC}
A continuación se presenta un circuito RC simple:

\begin{figure}[H]
    \centering
    \begin{circuitikz}[european]
        \draw (0,0) to[battery, l=$V$] (0,3)
              to[short] (2,3)
              to[R, l=$R$] (2,1.5)
              to[C, l=$C$] (2,0)
              to[short] (0,0);
    \end{circuitikz}
    \caption{Circuito RC simple}
    \label{fig:circuito_rc}
\end{figure}

La constante de tiempo $\tau$ del circuito RC es:
\[
\tau = R \cdot C
\]

\subsection{Circuito con Fuente de Corriente}

\begin{figure}[H]
    \centering
    \begin{circuitikz}[european]
        \draw (0,0) to[I, l=$I_s$] (0,3)
              to[short] (3,3)
              to[R, l=$R_1$] (3,0)
              to[short] (0,0);
        \draw (3,3) to[short] (5,3)
              to[L, l=$L$] (5,0)
              to[short] (3,0);
    \end{circuitikz}
    \caption{Circuito RL con fuente de corriente}
    \label{fig:circuito_rl}
\end{figure}

\section{Inserción de Imágenes PNG}

\subsection{Diagrama de Flujo}

\begin{figure}[H]
    \centering
    % \includegraphics[width=0.8\textwidth]{ejemplo-diagrama.png}
    \caption{Ejemplo de diagrama de flujo (imagen PNG)}
    \label{fig:diagrama_png}
\end{figure}

\textbf{Nota:} Reemplaza \texttt{ejemplo-diagrama.png} con el nombre de tu archivo de imagen.

\subsection{Múltiples Imágenes}

\begin{figure}[H]
    \centering
    \begin{subfigure}{0.45\textwidth}
        \centering
        % \includegraphics[width=\textwidth]{imagen1.png}
        \caption{Primera medición}
        \label{fig:medicion1}
    \end{subfigure}
    \hfill
    \begin{subfigure}{0.45\textwidth}
        \centering
        % \includegraphics[width=\textwidth]{imagen2.png}
        \caption{Segunda medición}
        \label{fig:medicion2}
    \end{subfigure}
    \caption{Comparación de mediciones experimentales}
    \label{fig:comparacion}
\end{figure}

\section{Tablas de Datos}

\subsection{Datos de Componentes}

\begin{table}[H]
    \centering
    \caption{Valores de componentes utilizados en el experimento}
    \label{tab:componentes}
    \begin{tabular}{l S[table-format=3.2] S[table-format=3.2] c}
        \toprule
        \textbf{Componente} & {\textbf{Valor Nominal}} & {\textbf{Valor Medido}} & \textbf{Tolerancia} \\
        \midrule
        Resistor R1  & \SI{1.00}{\kilo\ohm} & \SI{0.98}{\kilo\ohm} & $\pm$\SI{5}{\percent} \\
        Resistor R2  & \SI{2.20}{\kilo\ohm} & \SI{2.18}{\kilo\ohm} & $\pm$\SI{5}{\percent} \\
        Capacitor C1 & \SI{100}{\nano\farad} & \SI{102}{\nano\farad} & $\pm$\SI{10}{\percent} \\
        Capacitor C2 & \SI{470}{\micro\farad} & \SI{465}{\micro\farad} & $\pm$\SI{20}{\percent} \\
        Inductor L1  & \SI{10.0}{\milli\henry} & \SI{9.8}{\milli\henry} & $\pm$\SI{5}{\percent} \\
        \bottomrule
    \end{tabular}
\end{table}

\subsection{Resultados Experimentales}

\begin{table}[H]
    \centering
    \caption{Mediciones de voltaje y corriente en diferentes puntos del circuito}
    \label{tab:mediciones}
    \begin{tabular}{l S[table-format=1.2] S[table-format=1.2] S[table-format=1.2] S[table-format=1.2]}
        \toprule
        \textbf{Punto} & {\textbf{$V_{\text{teórico}}$ (\si{\volt})}} & {\textbf{$V_{\text{medido}}$ (\si{\volt})}} & {\textbf{$I_{\text{teórico}}$ (\si{\milli\ampere})}} & {\textbf{$I_{\text{medido}}$ (\si{\milli\ampere})}} \\
        \midrule
        Nodo A & 5.00 & 4.95 & 2.50 & 2.48 \\
        Nodo B & 3.30 & 3.28 & 1.65 & 1.63 \\
        Nodo C & 1.80 & 1.79 & 0.90 & 0.89 \\
        Nodo D & 0.00 & 0.01 & 0.00 & 0.01 \\
        \bottomrule
    \end{tabular}
\end{table}

\subsection{Análisis de Errores}

\begin{table}[H]
    \centering
    \caption{Análisis de error porcentual}
    \label{tab:errores}
    \begin{tabular}{l S[table-format=4.1] S[table-format=4.1] S[table-format=1.2]}
        \toprule
        \textbf{Variable} & {\textbf{Valor Teórico}} & {\textbf{Valor Experimental}} & {\textbf{Error (\si{\percent})}} \\
        \midrule
        Frecuencia de corte (\si{\hertz}) & 1591.5 & 1550.0 & 2.61 \\
        Ganancia máxima (\si{\deci\bel})   & 20.0   & 19.8   & 1.00 \\
        Ancho de banda (\si{\kilo\hertz})  & 10.0   & 9.8    & 2.00 \\
        \bottomrule
    \end{tabular}
\end{table}

\section{Fórmulas y Cálculos}

\subsection{Ecuaciones de Circuitos}

Para el circuito RC mostrado en la Figura \ref{fig:circuito_rc}, el voltaje en el capacitor es:

\[
V_C(t) = V_0 \left(1 - e^{-\frac{t}{RC}}\right)
\]

Donde:
\begin{itemize}
    \item $V_C(t)$ es el voltaje en el capacitor en el tiempo $t$
    \item $V_0$ es el voltaje de la fuente
    \item $R$ es la resistencia en ohmios
    \item $C$ es la capacitancia en faradios
\end{itemize}

\subsection{Impedancia Compleja}

La impedancia de un circuito RLC en serie es:

\[
Z = R + j\omega L + \frac{1}{j\omega C} = R + j\left(\omega L - \frac{1}{\omega C}\right)
\]

\section{Código de Simulación (Opcional)}

\begin{lstlisting}[caption={Código Python para simulación de circuito RC}, label=lst:simulacion]
import numpy as np
import matplotlib.pyplot as plt

# Parámetros de Circuito.
R = 1000  # Resistencia en ohmios
C = 1e-6  # Capacitancia en faradios
V0 = 5    # Voltaje de la fuente
tau = R * C  # Constante de tiempo

# Tiempo de Simulación.
t = np.linspace(0, 5*tau, 1000)

# Voltaje en el capacitor.
Vc = V0 * (1 - np.exp(-t/tau))

# Gráfica del resultado.
plt.figure(figsize=(10, 6))
plt.plot(t*1000, Vc, 'b-', linewidth=2, label='$V_C(t)$')
plt.xlabel('Tiempo (ms)')
plt.ylabel('Voltaje (V)')
plt.title('Respuesta de circuito RC')
plt.grid(True)
plt.legend()
plt.show()
\end{lstlisting}

\section{Conclusiones}
Esta plantilla proporciona una estructura organizada para documentar proyectos y experimentos relacionados con circuitos eléctricos. Se puede adaptar fácilmente para incluir:
\begin{itemize}
    \item Más tipos de circuitos usando \texttt{circuitikz}
    \item Gráficos de resultados experimentales
    \item Tablas adicionales de datos
    \item Cálculos matemáticos complejos
    \item Referencias bibliográficas
\end{itemize}

\section*{Recomendaciones}
\begin{enumerate}
    \item Mantener todas las imágenes en la misma carpeta del documento
    \item Utilizar nombres descriptivos para las figuras y tablas
    \item Verificar que los paquetes necesarios estén instalados
    \item Compilar con PDFLaTeX o XeLaTeX
\end{enumerate}

\end{document}